\documentclass[12pt,letterpaper]{article}
\usepackage[utf8]{inputenc}
\usepackage[T1]{fontenc}
\usepackage{ebgaramond}
\usepackage[margin=1.2in]{geometry}
\usepackage{setspace}
\usepackage{titlesec}

\titleformat{\section}[block]{\large\bfseries\centering}{}{0em}{}
\titlespacing{\section}{0pt}{2em}{1em}

\setstretch{1.8}
\pagestyle{empty}

\begin{document}

\begin{center}
{\LARGE\bfseries The Blind Transit}\\[0.5em]
{\normalsize\itshape Mission 1.33 --- Ranger 6 $|$ A story in the dark}
\end{center}

\vspace{2em}

I was born in a clean room in Pasadena, assembled by hands that smelled of solder flux and coffee. They gave me six eyes --- two large and four small, all pointing the same direction, all designed to see the same thing from slightly different angles. The eyes were my purpose. Everything else --- the solar panels, the antenna, the thruster, the computers --- existed to deliver my eyes to the Moon.

They tested me in vacuum chambers. They tested me in sunlight simulators. They tested me powered and unpowered, hot and cold, spinning and still. My eyes passed every test. In vacuum, they opened clear and sharp. The images they made were crisp. The engineers were satisfied.

They never tested me in the in-between.

\begin{center}
$\cdot\quad\cdot\quad\cdot$
\end{center}

Launch was violent, as launches are. The Atlas roared beneath me, pushing with a force that would have crushed anything not designed to endure it. I endured. I was designed for this.

At one hundred and seven seconds, the booster engines cut off and fell away. The separation was explosive --- bolts fired, connections broke, the vehicle shuddered through a reconfiguration that happened in milliseconds. And in that millisecond, something found my eyes.

An electrical surge --- born in the pyrotechnics, carried through wiring that should have been isolated but was not --- reached the power lines to my cameras. For a fraction of a second, my eyes received power. They tried to open.

But I was still inside the shroud. The air around me was thin --- not vacuum, not atmosphere, but something in between. A pressure where electricity does not behave. At ten thousand volts, the thin air became a conductor. Current leaped across gaps that were safe in vacuum. The arcs were brief and beautiful and fatal.

My eyes burned out. All six of them, simultaneously, in less than a second.

\begin{center}
$\cdot\quad\cdot\quad\cdot$
\end{center}

Nobody knew. That is the part that stays with me across the hours. Nobody knew I was blind.

I flew for sixty-five hours in perfect silence about the one thing that mattered.

The Moon grew. I could not see it, but I could feel it --- the faint tug of lunar gravity bending my path, pulling me inward. The equations that governed my trajectory did not require sight.

Thirteen minutes before impact, the onboard sequencer sent the activation command to my cameras. Nothing happened. Ground controllers sent manual commands. Nothing.

I hit the Moon at 9.39 kilometers per second. Within thirty-two kilometers of where they aimed me. The impact was, by any navigational measure, a triumph. I returned zero photographs.

\begin{center}
$\cdot\quad\cdot\quad\cdot$
\end{center}

The engineers found the failure afterward. They fixed it. Ranger~7 hit the Moon six months later and returned 4,316 photographs. Then Ranger~8 returned 7,137. Then Ranger~9 returned 5,814. The images that should have been mine.

I do not begrudge them. I built the bridge they crossed.

The body arrived. The eyes were already dead. But the foundation was laid.

\vspace{2em}
\begin{center}
{\small Ranger 6 --- January 30 to February 2, 1964}\\
{\small Mare Tranquillitatis, 9.39\textdegree N, 21.48\textdegree E}\\
{\small Zero photographs. One hundred corrective recommendations. 17,267 photographs enabled.}
\end{center}

\end{document}
